% Overleaf-ready thesis outline snippet.
% Assumes a thesis class with chapter-level numbering.
% If you want to paste this inside an existing chapter, replace \chapter with
% \section and shift the lower sectioning commands down one level.
% Recommended packages in the main document preamble:
% \usepackage{amsmath,amssymb}

\chapter{MS Thesis Topic}

This chapter follows the structure of the January 2026 thesis outline as closely as possible, but updates the text so it describes the final production system implemented in this repository rather than only the proposal-stage plan.

\section{Working Title}

Intervention-Aware Spatiotemporal Point Process Modeling for Predictive Multi-Robot Bird Deterrence in Vineyards

\section{Why This Topic Fits the Existing Framework}

The current production implementation already provides:

\begin{itemize}
\item health-weighted zone partitioning via a power diagram for UGV zones, with neighbor graph extraction,
\item a decentralized online spatiotemporal intensity estimator per robot (\texttt{OnlineSESTPP}) with self-excitation and cross-excitation near zone borders,
\item a task pipeline where detections generate immediate deterrence tasks and the predicted intensity field generates patrol and preventive-deterrence candidates via hotspot extraction,
\item an opt-in habituation-aware STL value function for predictive deterrence under per-cell, per-cue effectiveness loss,
\item a capability-aware task assigner that selects primary and secondary robots based on ETA, endurance, zone membership, health, load, and task value.
\end{itemize}

In the original outline, the key thesis extension was to make deterrence actions feed back into the event intensity model. The production system now implements that feedback explicitly through persistent intervention-inhibition channels in \texttt{SESTPP.py}, completion-triggered intervention updates in \texttt{DeterrentSystem.py}, and boundary sharing of intervention events between neighboring robots. The thesis therefore documents and evaluates an implemented intervention-aware predictive planning system rather than only proposing one.

\section{Core Research Question}

How effectively can an online spatiotemporal point process model that (i) predicts bird intrusion hotspots from distributed detections, (ii) explicitly represents the suppressive effect of deterrence actions, and (iii) accounts for cue habituation through an STL mission-value formulation drive multi-robot task selection so that future bird activity and value-weighted exposure are reduced more effectively than reactive, legacy predictive, or no-habituation baselines?

\section{Thesis Objectives}

\begin{enumerate}
\item Document and formalize the implemented intervention-aware online spatiotemporal intensity model by extending the self-exciting formulation with an inhibitory term driven by deterrence actions.
\item Integrate the model into the production task pipeline by scoring candidate patrol and deterrence locations using expected future intensity reduction or counterfactual STL robustness improvement relative to execution cost.
\item Evaluate whether a habituation-aware STL clause improves predictive action selection when repeated deterrence cues lose effectiveness.
\item Evaluate the final production system in simulation against strong baselines using metrics aligned with agricultural impact, energy, responsiveness, forecasting quality, and coordination overhead.
\end{enumerate}

\section{Problem Setting and Notation}

Let $\Omega \subset \mathbb{R}^2$ denote the vineyard area, or a discretized grid over that area. Robots are indexed by $r \in \mathcal{R}$ and each robot maintains a local model over its zone $\Omega_r$.

\begin{itemize}
\item Events (bird detections): $E_r = \{(x_i,t_i)\}$ are detections attributed to robot $r$ together with cross-shared neighbor events.
\item Interventions (deterrence actions): $U_r = \{(z_j,\tau_j,m_j)\}$ are completed deterrence actions executed by robot $r$ at location $z_j$ and time $\tau_j$ with deterrence mode or formation label $m_j$.
\item Habituation effectiveness: $\eta(z,m,t)$ is the current effectiveness of cue mode $m$ in cell $z$.
\item Value weighting: $w(x)$ denotes the spatial value surface over the vineyard, emphasizing rows, edges, or other high-value regions.
\end{itemize}

\section{Baseline Model: Online Self-Exciting Spatiotemporal Intensity}

The excitation component retained in the current \texttt{OnlineSESTPP} structure corresponds to a Hawkes-type intensity:

\begin{equation}
\lambda_r^{\mathrm{exc}}(x,t)
=
\mu_r(x)\,s(t)
+
\sum_{(x_i,t_i)\in E_r}
\alpha_{\mathrm{in}} K_{\sigma}(x-x_i)\exp\!\left(-\frac{t-t_i}{\omega}\right)
+
\sum_{(x_k,t_k)\in E_N(r)}
\alpha_{\mathrm{cross}} w_{r\leftarrow k} K_{\sigma}(x-x_k)\exp\!\left(-\frac{t-t_k}{\omega}\right),
\end{equation}

where $\mu_r(x)$ is a background rate, $s(t)$ is a diurnal multiplier, and $K_{\sigma}(\cdot)$ is a spatial Gaussian kernel.

\section{Implemented Extension: Intervention-Aware Self-Exciting and Self-Regulating Model}

The production system introduces an inhibitory, self-regulating component driven by deterrence actions:

\begin{equation}
\lambda_r^{\mathrm{IA}}(x,t)
=
\left[
\lambda_r^{\mathrm{exc}}(x,t)
-
\sum_{(z_j,\tau_j,m_j)\in U_r}
\beta(m_j,x,\tau_j)\,
K_{\sigma_u(m_j)}(x-z_j)\,
\exp\!\left(-\frac{t-\tau_j}{\omega_u(m_j)}\right)
\right]_+,
\end{equation}

where $[\cdot]_+ = \max(\cdot,0)$ enforces nonnegativity of intensity.

Interpretation:

\begin{itemize}
\item $\beta(\cdot)$ controls the immediate strength of deterrence,
\item $\omega_u(\cdot)$ controls how long deterrence effects persist,
\item $\sigma_u(\cdot)$ controls the spatial footprint of deterrence and can encode mode- or formation-dependent coverage.
\end{itemize}

In the production code, this is implemented through persistent per-mode inhibition channels whose sum forms the total inhibition mass used in the forecast field. The default production modes are \texttt{formation}, \texttt{laser}, and \texttt{biosonic}, each with its own suppression parameters.

\subsection{Implemented Habituation Model (Diminishing Returns)}

The current production runtime includes a per-cell, per-cue habituation state. Let $\eta(z,m,t) \in (0,1]$ denote the effectiveness of deterrence mode $m$ in cell $z$. Effectiveness recovers toward $1.0$ over time and decreases when the same cue is applied in the same cell.

The truth suppression contribution of a completed deterrence action is scaled by the effectiveness sampled at completion time:

\begin{equation}
s_{j}(x,t)
=
\eta(z_j,m_j,\tau_j)\,
\beta_0(m_j)\,
K_{\sigma_u(m_j)}(x-z_j)\,
\exp\!\left(-\frac{t-\tau_j}{\omega_u(m_j)}\right).
\end{equation}

Thus repeated cue use can reduce realized suppression in the simulated world. The control condition sets habituation off, or equivalently sets the habituation drop parameter to zero, which keeps $\eta=1$ at application time.


\subsection{Habituation-Aware STL Predictive Value}

The production planner can replace the legacy $\Delta J$ predictive value with a counterfactual STL robustness improvement:

\begin{equation}
U(a,r)
=
\rho(\Phi_r,\xi^{a}) - \rho(\Phi_r,\xi^{0}),
\end{equation}

where $\xi^{a}$ is the predicted local trace with candidate action $a$, $\xi^{0}$ is the trace without that action, and $\Phi_r$ is a robot-local mission formula. The production formula uses exposure, coverage, and optional habituation clauses:

\begin{equation}
\Phi_r = \varphi_{\mathrm{exp}} \wedge \varphi_{\mathrm{cov}} \wedge \varphi_{\mathrm{hab}}.
\end{equation}

The no-habituation STL baseline uses only exposure and coverage. The full STL proposal adds $\varphi_{\mathrm{hab}}$, which penalizes candidates that reuse a cue whose current effectiveness is low. The resulting $U(a,r)$ is copied into the existing utility and score fields, so the dispatcher is unchanged.
\section{Planning: Scoring Actions by Expected Future Reduction}

At each replan cycle, the production system generates candidate targets from the current predicted intensity field, then scores each candidate deterrence action by its expected reduction in future intensity over a horizon $T_H$.

Let a candidate action be $a=(z,t,m)$ executed at current time $t$. A continuous reduction objective consistent with the original outline is:

\begin{equation}
\Delta J(a)
=
\int_t^{t+T_H}\int_{\Omega_r}
w(x)\left(
\lambda_r^{\mathrm{IA}}(x,\tau)
-
\lambda_r^{\mathrm{IA},a}(x,\tau)
\right)\,dx\,d\tau
- c_E(a),
\end{equation}

where $w(x)$ is the value weighting over the vineyard and $c_E(a)$ is an energy or time cost.

A computationally efficient approximation follows from the exponential temporal kernel:

\begin{equation}
\int_t^{t+T_H}
\exp\!\left(-\frac{\tau-t}{\omega_u}\right)\,d\tau
=
\omega_u\left(1-\exp\!\left(-\frac{T_H}{\omega_u}\right)\right).
\end{equation}

The production implementation uses this same structure in a local grid-patch approximation in \texttt{planner\_task\_estimation.py}. For each candidate action, the runtime computes a capped counterfactual exposure reduction exported as \texttt{predicted\_deltaJ}, then forms the operational utility:

\begin{align}
U_{\mathrm{det}}(a) &= \Delta J(a) - c_{\eta}(a), \\
\rho_{\Delta J}(a) &= \frac{\Delta J(a)}{\max(c_{\eta}(a),\varepsilon)}.
\end{align}

Patrol candidates are scored with an analogous horizon-weighted excess-field benefit minus ETA cost, so both predictive branches compete inside the same dispatch pipeline.

\section{Decentralized Neighbor Sharing With Interventions}

The existing boundary-event mechanism is extended in production to share intervention events near zone boundaries in addition to detection events:

\begin{itemize}
\item If a deterrence action completes within a distance $d_b$ of the zone boundary, the runtime broadcasts a compact message containing the action location, time, mode, strength, footprint, and decay parameters.
\item Neighboring robots add that event into their local inhibitory term as cross-zone intervention feedback, improving prediction consistency near borders.
\item Communication remains efficient because it is local, event-triggered, and filtered by debounce interval, spatial quantization, and minimum message weight.
\end{itemize}

\section{Implementation in the Existing Codebase}

\subsection{Model Extension (\texttt{SESTPP.py})}

The production \texttt{SESTPP.py} implementation keeps a second family of grid accumulators for inhibition mass, analogous to trigger mass:

\begin{itemize}
\item per-mode inhibition channels are updated by \texttt{add\_intervention\_event(...)} using spatial stamping similar to the excitation stamp,
\item in \texttt{advance\_time(dt)}, trigger mass decays with $\exp(-dt/\omega)$ and each inhibition channel decays with its own $\exp(-dt/\omega_u)$,
\item the runtime forecast is computed as the clipped background plus trigger mass minus the summed inhibition mass.
\end{itemize}

\subsection{Closed-Loop Integration (\texttt{DeterrentSystem.py})}

When a deterring task completes, the production runtime:

\begin{itemize}
\item calls the robot-local intervention update so the completed action is written back into the robot model,
\item adds the same action to the recent intervention history used by the ground-truth suppression process,
\item broadcasts intervention boundary events to neighboring robots when the completed action occurs near the zone boundary,
\item records the event in structured runtime snapshots and communication metrics.
\end{itemize}

\subsection{Action Scoring (\texttt{TaskGenerator.py} and \texttt{planner\_task\_estimation.py})}

The production planner augments the hotspot-selection path with model-scored preventive actions:

\begin{itemize}
\item candidate points are generated from forecast hotspots, border hotspots, and clustered recent detections,
\item for each candidate deterrence action, the runtime computes $\Delta J(a)$ with the counterfactual reduction estimator,
\item patrol candidates are also scored over the same planning horizon,
\item the best-scoring tasks are passed forward to assignment and dispatch subject to persistence, ETA, queue-cap, and conflict protections.
\end{itemize}

\subsection{Dispatch, Motion, and Tracking}

Beyond the original outline, the final production system also includes:

\begin{itemize}
\item dispatch-time admission logic in \texttt{planner\_dispatch.py},
\item row-constrained motion execution and headland lane switching in \texttt{DeterrentSystem.py},
\item structured stage outputs and runtime snapshots in \texttt{system\_structure.py},
\item telemetry and tracking export for experiments, diagnostics, and live demos.
\end{itemize}

\section{Experimental Design}

\subsection{Simulation Scenarios}

The production evaluation uses a ground-truth generative process for bird events so interventions can have measurable impact:

\begin{itemize}
\item base rate plus self-excitation to generate clustered bird activity,
\item row and edge value weighting through $w(x)$,
\item a detection model with false negatives and range limits consistent with sensor realism,
\item intervention effects applied to the ground-truth intensity to represent actual deterrence impact.
\end{itemize}

\subsection{Baselines}

At minimum:

\begin{enumerate}
\item Reactive: respond only to detections, with no predictive patrol.
\item Prediction-only: use the self-exciting forecast model for hotspot generation, but do not let deterrence actions suppress predicted risk.
\item Proposed / production intervention-aware system: use intervention-aware forecasting together with patrol and preventive action scoring.
\end{enumerate}

Optional ablations remain useful for the final thesis:

\begin{itemize}
\item no neighbor sharing vs. sharing detections only vs. sharing detections plus interventions,
\item no intervention feedback vs. intervention feedback,
\item heuristic preventive admission vs. later selective-admission variants,
\item no habituation vs. habituation-aware suppression and STL scoring.
\end{itemize}

\subsection{Metrics}

Report mean and variance over multiple randomized runs. The primary exposure metric remains:

\begin{equation}
J_{\mathrm{exp}}
=
\int_0^T \int_{\Omega}
w(x)\,\lambda_{\mathrm{true}}(x,t)\,dx\,dt.
\end{equation}

Additional metrics include:

\begin{itemize}
\item mean response time from event onset to first deterrence arrival,
\item total energy or travel distance per robot, with UGV and UAV accounting separated when needed,
\item task efficiency, such as completed tasks per unit distance or exposure per completed task,
\item communication overhead, including boundary-message counts and estimated bytes sent,
\item forecast quality, including recall and precision at top-$k$ hotspots,
\item realized suppression metrics such as truth suppression rate and birds deterred percentage,
\item habituation metrics such as cue effectiveness at application and cue-variety index,
\item STL robustness metrics for the exposure, coverage, and habituation clauses.
\end{itemize}

\section{Implemented Contributions}

\begin{itemize}
\item A practical online intervention-aware spatiotemporal point process model suitable for real-time multi-robot operation.
\item Planning objectives that select actions based on predicted future reduction and, in the STL mode, counterfactual mission-robustness improvement.
\item A decentralized coordination mechanism that shares both detections and intervention effects locally to reduce border artifacts.
\item A closed-loop simulation and evaluation framework that measures value-weighted exposure, responsiveness, deterrence effectiveness, habituation, STL robustness, forecast quality, and communication cost.
\end{itemize}

\section{Implemented Deliverables and Stretch Goals}

Implemented deliverables:

\begin{itemize}
\item intervention-aware forecast dynamics integrated into the production simulation loop,
\item action scoring based on counterfactual future reduction and habituation-aware STL robustness together with patrol scoring and dispatch integration,
\item a full experimental workflow with baselines, calibration, diagnostics, and summary metrics.
\end{itemize}

Remaining extensions, if needed for future work:

\begin{itemize}
\item add online parameter adaptation for $\beta_0(m)$ and $\omega_u(m)$ using likelihood-based updates,
\item study formation-dependent footprints by mapping deterrence mode to $\sigma_u(m)$ and comparing footprints empirically.
\end{itemize}

\section{Proposed Thesis Chapter Outline}

\begin{enumerate}
\item Introduction and problem motivation.
\item Related work: persistent monitoring, hotspot prediction, spatiotemporal point processes, intervention modeling, and multi-robot deterrence.
\item System overview and existing architecture.
\item Intervention-aware point process model: formulation and online update.
\item Planning and task selection using expected future reduction.
\item Experimental setup and baselines.
\item Results, ablations, habituation-aware STL evidence, and discussion.
\item Conclusions and future work.
\end{enumerate}
